% =========================================================
% Summary Tables
% =========================================================

\subsection*{Summary Tables}

<<<<<<< HEAD
\begin{toolkitbox}{Predicate Logic Summary Toolkit --- Quick Reference}
\small
\begin{tabular}{p{0.30\linewidth} p{0.38\linewidth} p{0.22\linewidth}}
\toprule
\textbf{Reference area} & \textbf{Use it to check} & \textbf{Main danger} \\
\midrule
Syntax objects & What kind of expression has been formed & Treating open formulas as sentences \\
Semantic objects & What must be supplied to evaluate a formula & Forgetting assignments \\
Translation patterns & How English quantifier phrases become formulas & Using the wrong connective \\
Proof cues & Which rule matches the goal or premise & Illegal witnesses \\
\bottomrule
\end{tabular}
\end{toolkitbox}

\begin{remark*}[Orientation]
These tables collect the local working vocabulary for predicate logic. They
are meant for quick lookup while translating statements and checking proofs.
When a table gives a schema, variables and predicate letters are placeholders;
the surrounding domain and intended predicates must still be supplied.
\end{remark*}

\subsection*{Syntax Objects}

\begin{center}
\renewcommand{\arraystretch}{1.25}
\begin{tabular}{p{0.22\linewidth} p{0.35\linewidth} p{0.33\linewidth}}
\toprule
\textbf{Object} & \textbf{Role} & \textbf{Example} \\
\midrule
Variable & Placeholder that may be bound or free & $x,y,z$ \\
Constant symbol & Name-like symbol of the language & $0$, $c$ \\
Function symbol & Symbol forming terms from terms & $f(x)$, $+(x,y)$ \\
Predicate symbol & Relation symbol forming atomic formulas & $P(x)$, $R(x,y)$ \\
Term & Object-denoting expression & $x$, $c$, $f(x)$ \\
Atomic formula & Predicate applied to terms, or equality of terms & $R(x,c)$, $t=s$ \\
Formula & Expression built using connectives and quantifiers & $\forall x\,(P(x)\rightarrow Q(x))$ \\
Sentence & Formula with no free variables & $\exists x\,P(x)$ \\
Open formula & Formula with at least one free variable & $P(x)\wedge Q(y)$ \\
\bottomrule
\end{tabular}
\end{center}

\subsection*{Semantics Objects}

\begin{center}
\renewcommand{\arraystretch}{1.25}
\begin{tabular}{p{0.24\linewidth} p{0.40\linewidth} p{0.26\linewidth}}
\toprule
\textbf{Object} & \textbf{Role} & \textbf{Notation} \\
\midrule
Structure & Supplies a domain and interpretations for symbols & $\mathcal M$ \\
Variable assignment & Sends variables to objects in the domain & $s$ \\
Modified assignment & Changes one variable value and leaves the rest fixed & $s[x\mapsto a]$ \\
Term evaluation & Value of a term in a structure under an assignment & $t^{\mathcal M}[s]$ \\
Satisfaction & Formula holds in a structure under an assignment & $\mathcal M,s\models\varphi$ \\
Truth in a structure & Sentence holds in a structure & $\mathcal M\models\sigma$ \\
Model of a theory & Structure satisfying every sentence in a theory & $\mathcal M\models T$ \\
\bottomrule
\end{tabular}
\end{center}

\begin{remark*}[Interpretation]
Syntax asks whether an expression is well formed. Semantics asks whether that
well-formed expression is true relative to a structure, and, for open formulas,
relative to an assignment. Sentences are special because their truth does not
depend on a variable assignment.
\end{remark*}

\subsection*{Quantifier Translations}

\begin{center}
\renewcommand{\arraystretch}{1.3}
\begin{tabular}{p{0.30\linewidth} p{0.38\linewidth} p{0.22\linewidth}}
\toprule
\textbf{English form} & \textbf{Logical form} & \textbf{Cue} \\
\midrule
All $A$ are $B$ & $\forall x\,(A(x)\rightarrow B(x))$ & Universal uses implication \\
No $A$ are $B$ & $\forall x\,(A(x)\rightarrow \neg B(x))$ & Universal exclusion \\
Some $A$ are $B$ & $\exists x\,(A(x)\wedge B(x))$ & Existential uses conjunction \\
Some $A$ are not $B$ & $\exists x\,(A(x)\wedge \neg B(x))$ & Counterexample form \\
Every $A$ has a $B$ & $\forall x\,(A(x)\rightarrow \exists y\,(B(y)\wedge R(x,y)))$ & Witness may depend on $x$ \\
There is one $B$ for every $A$ & $\exists y\,(B(y)\wedge \forall x\,(A(x)\rightarrow R(x,y)))$ & Uniform witness \\
\bottomrule
\end{tabular}
\end{center}

\subsection*{Quantifier Negation Laws}

\begin{center}
\renewcommand{\arraystretch}{1.35}
\begin{tabular}{p{0.36\linewidth} p{0.36\linewidth} p{0.18\linewidth}}
\toprule
\textbf{Negated form} & \textbf{Equivalent form} & \textbf{Reading} \\
\midrule
$\neg\forall x\,\varphi$ & $\exists x\,\neg\varphi$ & A counterexample exists \\
$\neg\exists x\,\varphi$ & $\forall x\,\neg\varphi$ & No witness exists \\
=======
\begin{remark*}[English phrase to logical pattern]
\renewcommand{\arraystretch}{1.35}
\begin{tabular}{p{0.32\linewidth} p{0.58\linewidth}}
\toprule
\textbf{English phrase} & \textbf{Logical pattern} \\
\midrule
All \(A\) are \(B\) & \(\forall x\,(A(x)\to B(x))\) \\
Some \(A\) are \(B\) & \(\exists x\,(A(x)\land B(x))\) \\
No \(A\) are \(B\) & \(\forall x\,(A(x)\to \neg B(x))\) \\
Some \(A\) are not \(B\) & \(\exists x\,(A(x)\land \neg B(x))\) \\
Every \(A\) has a \(B\) & \(\forall x\,(A(x)\to \exists y\,(B(y)\land R(x,y)))\) \\
There is a \(B\) for every \(A\) & \(\exists y\,(B(y)\land \forall x\,(A(x)\to R(x,y)))\) \\
>>>>>>> 591349ed2a1982198e862953b73959c91a638d59
\bottomrule
\end{tabular}
\end{remark*}

<<<<<<< HEAD
\subsection*{Proof-Rule Cues}

\begin{center}
\renewcommand{\arraystretch}{1.25}
\begin{tabular}{p{0.25\linewidth} p{0.35\linewidth} p{0.30\linewidth}}
\toprule
\textbf{Rule} & \textbf{Use when} & \textbf{Side condition} \\
\midrule
Universal instantiation & You have $\forall x\,\varphi$ and need an instance & Term must be free for $x$ \\
Universal generalization & You have proved the arbitrary case & No dependence on special assumptions \\
Existential instantiation & You have $\exists x\,\varphi$ and need a witness name & Use a fresh temporary constant \\
Existential generalization & You have one instance and need existence & Replace the instance by a quantified variable \\
Equality substitution & You know $s=t$ and need to replace equals & Substitute only in valid positions \\
\bottomrule
\end{tabular}
\end{center}

\subsection*{Common Proof Strategies}

\begin{center}
\renewcommand{\arraystretch}{1.25}
\begin{tabular}{p{0.25\linewidth} p{0.34\linewidth} p{0.31\linewidth}}
\toprule
\textbf{Goal} & \textbf{Strategy} & \textbf{Check} \\
\midrule
Prove universal & Let an arbitrary object be given, then prove the instance & The object must remain arbitrary \\
Prove existential & Construct or identify one witness & Verify all required properties \\
Disprove universal & Produce a counterexample & Show the object satisfies the domain condition \\
Disprove existential & Prove every possible witness fails & Usually rewrite as a universal negation \\
Prove implication & Assume the antecedent, prove the consequent & Discharge the assumption properly \\
Prove equivalence & Prove both implications & Track each direction separately \\
\bottomrule
\end{tabular}
\end{center}

\subsection*{Common Fallacies}

\begin{center}
\renewcommand{\arraystretch}{1.25}
\begin{tabular}{p{0.28\linewidth} p{0.34\linewidth} p{0.28\linewidth}}
\toprule
\textbf{Fallacy} & \textbf{Bad form} & \textbf{Correction} \\
\midrule
Universal/existential confusion & $\forall x\,P(x)$ for ``some $P$'' & Use $\exists x\,P(x)$ \\
All with conjunction & $\forall x\,(A(x)\wedge B(x))$ & $\forall x\,(A(x)\rightarrow B(x))$ \\
Some with implication & $\exists x\,(A(x)\rightarrow B(x))$ & $\exists x\,(A(x)\wedge B(x))$ \\
Invalid quantifier reversal & Swap $\forall x\exists y$ with $\exists y\forall x$ & Preserve dependency order \\
Free variable left open & $P(x)$ as a sentence & Bind $x$ or state the assignment context \\
Variable capture & Substitute into a binding scope blindly & Rename bound variables first \\
Illicit generalization & Generalize from a special object & Prove the arbitrary case \\
Illicit instantiation & Reuse an old constant as a witness & Introduce a fresh witness \\
Affirming the consequent & $P\rightarrow Q$, $Q$, therefore $P$ & Use modus ponens or contrapositive \\
Denying the antecedent & $P\rightarrow Q$, $\neg P$, therefore $\neg Q$ & Use modus tollens only with $\neg Q$ \\
\bottomrule
\end{tabular}
\end{center}

\begin{example}[Reading a quantifier pattern]\label{ex:predicate-logic-summary-quantifier-pattern}
The formulas
\[
\forall x\exists y\,R(x,y)
\qquad\text{and}\qquad
\exists y\forall x\,R(x,y)
\]
have the same symbols but different dependency structure. In the first formula
the witness $y$ may vary with $x$. In the second formula one witness must work
for every $x$.
\end{example}
=======
\begin{remark*}[Quantifier negation laws]
\[
\begin{array}{rcl}
\neg\forall x\,\varphi &\Longleftrightarrow& \exists x\,\neg\varphi,\\
\neg\exists x\,\varphi &\Longleftrightarrow& \forall x\,\neg\varphi,\\
\neg\forall x\exists y\,\varphi &\Longleftrightarrow&
  \exists x\forall y\,\neg\varphi,\\
\neg\exists x\forall y\,\varphi &\Longleftrightarrow&
  \forall x\exists y\,\neg\varphi.
\end{array}
\]
\end{remark*}

\begin{remark*}[Categorical form translations]
\renewcommand{\arraystretch}{1.35}
\begin{tabular}{p{0.12\linewidth} p{0.28\linewidth} p{0.48\linewidth}}
\toprule
\textbf{Form} & \textbf{Name} & \textbf{Translation} \\
\midrule
A & Universal affirmative & \(\forall x\,(A(x)\to B(x))\) \\
E & Universal negative & \(\forall x\,(A(x)\to \neg B(x))\) \\
I & Particular affirmative & \(\exists x\,(A(x)\land B(x))\) \\
O & Particular negative & \(\exists x\,(A(x)\land \neg B(x))\) \\
\bottomrule
\end{tabular}
\end{remark*}

\begin{remark*}[Proof strategy cues]
\renewcommand{\arraystretch}{1.35}
\begin{tabular}{p{0.28\linewidth} p{0.58\linewidth}}
\toprule
\textbf{Goal or premise} & \textbf{Useful cue} \\
\midrule
Prove \(\forall x\,\varphi(x)\) & Let \(x\) be arbitrary, then prove \(\varphi(x)\). \\
Use \(\forall x\,\varphi(x)\) & Instantiate it with a term that is free for \(x\). \\
Prove \(\exists x\,\varphi(x)\) & Produce a witness and verify \(\varphi\). \\
Use \(\exists x\,\varphi(x)\) & Introduce a fresh witness and keep its scope controlled. \\
Disprove \(\forall x\,\varphi(x)\) & Find a counterexample satisfying \(\neg\varphi\). \\
Disprove an implication & Make the hypothesis true and the conclusion false. \\
\bottomrule
\end{tabular}
\end{remark*}
>>>>>>> 591349ed2a1982198e862953b73959c91a638d59
